% #####################################################################
\chapter{\texorpdfstring{Note on the Use of AI}{Note on the Use of AI}}\label{chap:nota-ia}
% #####################################################################
% Versão curta (2026-11-15). A versão longa ficou fora do texto.

\textcolor{red}{Este trabalho usou o Claude Code, um assistente de programação da Anthropic que roda no terminal e lê e edita os arquivos do repositório. O uso foi de maio a novembro de 2026.}

\textcolor{red}{O assistente ajudou a resumir as referências, a escrever e depurar o simulador e os scripts dos testes em ponto fixo, o código HLS e os scripts de síntese, implementação e potência. Ele também escreveu os scripts das figuras e tabelas e rascunhos de seções em português, que a autora reescreveu e traduziu.}

\textcolor{red}{A autora escolheu o tema, o escopo e o artigo base. Ela definiu os critérios de escolha dos formatos, a FPGA alvo e as arquiteturas testadas, fez as rodadas no servidor e decidiu o que cada figura conclui. A versão final do texto é dela.}

\textcolor{red}{O assistente errou em alguns pontos, como propor referências de memória e calcular um tempo com a média errada de ciclos. A autora corrigiu esses erros na revisão. Por isso, só entram citações de documentos que a autora tem.}

% CONFERIR: a autora confirma se essa frase vale para todo o texto.
\textcolor{red}{Os números do texto vêm dos relatórios e dos registros das ferramentas, e não da resposta do assistente.}
